\section{Introduction}\label{sec:introduction}

Interactions move particles out of the condensate even in a Bose gas at
zero temperature. For a dilute three-dimensional gas, Bogoliubov theory
predicts both the number of these depleted particles and their distribution
in momentum. Establishing these predictions for an exact ground state
requires information at arbitrarily long wavelengths: their energy cost
tends to zero as the volume grows. This paper determines the leading
momentum distribution for hard spheres, with the ordinary thermodynamic
limit taken before the dilute limit.

\subsection{The model and the result}

Fix $a>0$ and write $\Lambda_L=(\R/L\Z)^3$ for the flat torus of side
$L>2a$, with torus distance $d_L$. The allowed configuration space is
\[
 \Omega_{N,L,a}=\{(x_1,\ldots,x_N)\in\Lambda_L^N:
                    d_L(x_i,x_j)>a\text{ for }i<j\}.
\]
The hard-sphere Hamiltonian $H^{\mathrm{hc}}_{N,L,a}$ is the operator
associated with the quadratic form
\[
 q[\Psi]=\sum_{i=1}^N\int_{\Omega_{N,L,a}}|\nabla_i\Psi|^2,
 \qquad
 \mathcal D(q)=H_0^1(\Omega_{N,L,a})\cap L^2_{\mathrm{sym}}(\Omega_{N,L,a}).
\]
Here $H_0^1$ is the $H^1$ closure of smooth functions compactly supported
in the allowed configuration set; the ambient torus is already periodic.
We use units $\hbar^2/(2m)=1$. The parameter
$a$ is both the exclusion distance between centers and the scattering
length: the exterior zero-energy scattering solution is $1-a/|x|$.
The form domain is nontrivial along every sequence with
$N/L^3\to\rho<1/(8a^3)$ and $N,L\to\infty$: a cubic grid of spacing
at least $2a$ eventually provides $N$ sites, and small neighborhoods
of these sites give an allowed open configuration set. Zero extension
and compactness on the full torus then give compact resolvent.

Let $P_{N,L,a}$ be the projection onto the lowest bosonic eigenspace,
and let $\mathcal G_{N,L,a}$ consist of the positive trace-one operators
$\Gamma$ satisfying $\Gamma=P_{N,L,a}\Gamma P_{N,L,a}$. After extension
by zero to the full periodic configuration space, put
\[
 \gamma^{(1)}_\Gamma=N\Tr_{2,\ldots,N}\Gamma,
 \qquad u_{k,L}(x)=L^{-3/2}e^{ik\cdot x}\quad
                    (k\in(2\pi/L)\Z^3),
\]
and define
\[
 n_{\Gamma,L}(k)=\langle u_{k,L},\gamma^{(1)}_\Gamma u_{k,L}\rangle,
 \qquad
 B_{N,L,a}(\Gamma)=\frac{\langle u_{0,L},\gamma^{(1)}_\Gamma u_{0,L}\rangle}{N}.
\]
Thus $B_{N,L,a}$ is the occupation fraction of the specified constant
orbital, and $1-B_{N,L,a}$ is its depletion.
The total occupation is $\sum_k n_{\Gamma,L}(k)=N$.
For $t\in\R^3\setminus\{0\}$, set
\begin{equation}\label{intro:profile}
 g(t)=\frac{|t|^2+1}{2\sqrt{|t|^4+2|t|^2}}-\frac12,
 \qquad
 \nu_{\mathrm{Bog}}(\dd t)
       =\frac{(8\pi)^{3/2}}{(2\pi)^3}g(t)\,\dd t.
\end{equation}
This is a finite positive measure: $g(t)=O(|t|^{-1})$ near zero
and $g(t)=O(|t|^{-4})$ at infinity. Its total mass is $8/(3\sqrt\pi)$.

\begin{theorem}\label{thm:main}
For every $a>0$, every bounded continuous real function $f$ on $\R^3$,
and every $\eps>0$, there is $\rho_0(a,f,\eps)>0$ with the following
property. Fix $0<\rho<\rho_0$ and put $\eta=\rho a^3$. For every
sequence $N_j,L_j\longrightarrow\infty$ with
$N_j/L_j^3\longrightarrow\rho$,
\begin{equation}\label{intro:main-limit}
 \limsup_{j\to\infty}\ \sup_{\Gamma\in\mathcal G_{N_j,L_j,a}}
 \left|
 \frac{1}{N_j\sqrt\eta}
 \sum_{\substack{k\in(2\pi/L_j)\Z^3\\k\ne0}}
 f\!\left(\frac{k}{\sqrt{8\pi\rho a}}\right)n_{\Gamma,L_j}(k)
       -\int_{\R^3}f\,\dd\nu_{\mathrm{Bog}}
 \right|\le\eps.
\end{equation}
\end{theorem}

The density is fixed throughout each thermodynamic limit. In particular,
the theorem controls every thermodynamic accumulation value of the
occupation, without requiring a unique occupation limit at fixed
positive $\eta$. Its uniformity includes all normalized, possibly
complex ground vectors and the normalized projection onto the entire
ground eigenspace.

\begin{corollary}\label{cor:depletion}
In the same iterated limit, uniformly over all ground states,
\[
 1-B_{N,L,a}(\Gamma)
   =\frac8{3\sqrt\pi}\sqrt{\rho a^3}+o(\sqrt{\rho a^3}).
\]
\end{corollary}
This follows by taking $f=1$. The theorem also determines how the
depleted particles are distributed on the momentum scale
$\sqrt{8\pi\rho a}$, without assuming translation invariance of an
individual ground state.

\subsection{Historical context and the sources of the method}

Bogoliubov's description of a weakly interacting Bose gas replaces its
low-energy excitations by quasiparticles obtained from a linear canonical
transformation~\cite{Bogoliubov}. The quasiparticle vacuum already has
particles in nonzero momentum modes. Lee, Huang, and Yang developed the
hard-sphere pseudopotential calculation far enough to determine both the
second term of the energy and the leading depletion~\cite{LHY}. Thus the
predicted momentum profile and the coefficient $8/(3\sqrt\pi)$ describe
the many-body state, as well as the fluctuations responsible for the
energy correction. The distinction matters mathematically: matching
the energy does not by itself control the occupation of modes whose
energy tends to zero with the inverse size of the system.

The rigorous energy theory began with Dyson's bounds for hard
spheres~\cite{Dyson}; Lieb and Yngvason later obtained the matching
leading lower bound~\cite{LY}. At the Lee--Huang--Yang order, Yau and
Yin proved the upper bound for smooth repulsive short-range
potentials~\cite{YY}. Basti, Cenatiempo, and Schlein subsequently gave
a proof for nonnegative radial compactly supported $L^3$ potentials
with an improved error~\cite{BCS2021}. Fournais and Solovej established
the lower bound, first for integrable interactions and then for a class
including hard cores~\cite{FS2020,FS2023}. For hard spheres, Basti,
Cenatiempo, Giuliani, Olgiati, Pasqualetti, and Schlein bounded the
correction by a constant times the Lee--Huang--Yang scale~\cite{BCGOPS};
Basti, Brooks, Cenatiempo, Olgiati, and Schlein obtained the sharp
coefficient~\cite{BBCOS}. Their subsequent work extends this upper-bound
method to measurable even nonnegative compactly supported potentials,
including hard cores~\cite{BBCOSgeneral}.
The hard-sphere upper bound supplies the energy comparison used here.

Rigorous information about the state has developed in several different
scaling limits. Lieb and Seiringer proved complete condensation for
trapped gases in the Gross--Pitaevskii limit, allowing hard cores; their
result includes trace-norm convergence of the normalized one-particle
density matrix and the corresponding leading momentum-density
limit~\cite{LS2002}. For a periodic unit box with scattering length of
order $N^{-1}$, Boccato, Brennecke, Cenatiempo, and Schlein determined
the next-order depletion for nonnegative radial compactly supported
$L^3(\R^3)$ potentials~\cite[Equation~(1.14) and Appendix~A]{BBCS}.
The associated discrete momentum occupations are expressed as a
trace-norm expansion of the one-particle density matrix by Brennecke,
Lee, and Nam, who also establish its positive-temperature
analogue~\cite[Equations~(6)--(8) and Theorem~1.1]{BLN}.
The second-order density-matrix expansion resolves the individual
discrete occupations as well as their sum. Its relation between volume
and scattering length differs from the thermodynamic limit at fixed
density considered here.

Condensation on boxes growing with the dilute parameter was proved by
Fournais~\cite{Fournais2021}, and the propagation of such estimates
to larger scales has been developed further by Junge~\cite{Junge2026}.
These results keep the box size tied to the dilution. In the ordinary
thermodynamic limit, the one-particle gap vanishes at each fixed
positive density, so an additional estimate is needed for variation
across the whole torus. A different comparison is provided by the
``simple equation'' approximation to the Bose gas: Carlen, Jauslin,
and Lieb derived its leading depletion coefficient~\cite[Theorem~1.6]{CJL},
and Jauslin derived its Bogoliubov momentum profile on the inverse
healing-length scale~\cite[Theorem~3]{Jauslin2024}.
Those are results for an approximate correlation equation. The present
occupation measures are those of the exact many-body hard-sphere
ground states, and their convergence is weak convergence of finite
measures on the rescaled momentum space.

Our local energy analysis builds on the lower-bound methods of
Fournais and Solovej~\cite{FS2020,FS2023} and the Neumann-box
localization of Haberberger, Hainzl, Nam, Seiringer, and
Triay~\cite{HHNST}. Fournais, Junge, Girardot, Morin, Olivieri, and
Triay combined these approaches in their treatment of strong
interactions at low temperature~\cite{FJGMOT}. We use intermediate
operator estimates from that work, retaining a positive term that
measures Bogoliubov excitations, before taking an observable expectation.
This is a use of their operator analysis beyond its scalar free-energy
conclusion. The excitation-number localization is related to the
large-matrix localization of Lieb and Solovej~\cite[Appendix~A]{LS2001},
and the coherent-state treatment of the condensate mode uses the
$c$-number substitution of Lieb, Seiringer, and Yngvason~\cite{LSY2005}.

The long-distance comparison uses a stationary ground-state path
representation, a form of Doob's $h$-transform~\cite{Doob}. The spatial
count estimates and exact conditional killed-Brownian-path laws used
here are supplied by the companion paper~\cite{H}. The companion
establishes a positive constant-orbital occupation uniformly over the
hard-sphere ground space. Here the comparison of particle-deletion
measures must be accurate below the depletion scale. Its additional
quantitative ingredients are a local fourth-moment comparison and a
distant second-moment comparison. These compare the deletion measures
obtained by averaging the deleted position over different cells.

The route construction adapts the unpredictable-path method of
Benjamini, Pemantle, and Peres, who constructed paths in three dimensions
with exponential intersection tails~\cite{BPP}. Our comparison also
requires encounters to become rare after the segments near the common
endpoint are removed, and it must accommodate finite defective regions.
Detours use Kesten's exterior-boundary connectivity theorem in the
formulation of Tim\'ar~\cite[Theorem~4]{Timar}. These geometric inputs
are combined with two-copy likelihood estimates that retain the exact
conditional normalization; the
intersection theorem alone does not give the deletion-measure comparison.

\subsection{From local excitations to the momentum profile}

The proof separates excitations inside moderately large boxes from
variation between those boxes. For a partition into cubes of side $\ell$,
let $P_\ell$ be the one-particle orthogonal projection onto functions
constant on each cube, and let $P_0$ be projection onto the single global
constant. Since $P_0\le P_\ell$, every normalized ground vector satisfies
the exact identity
\begin{equation}\label{intro:split}
 1-B_{N,L,a}(|\Psi\rangle\langle\Psi|)
 =\langle\Psi,(1-P_\ell)_1\Psi\rangle
   +\|(P_\ell-P_0)_1\Psi\|^2.
\end{equation}
The subscript $1$ means that the projection acts on the first particle.
Bosonic symmetry gives the same identity with the average over all particles.

The first term can be analyzed through the energy, provided the lower
bound retains a positive operator measuring Bogoliubov excitations.
The energy boxes contain several of the projection cells of side $\ell$;
those cells will in turn be subdivided into small cells for the spatial
comparison.
An energy expansion alone does not identify this observable. We use
the hard-sphere upper bound of Basti, Brooks, Cenatiempo, Olgiati,
and Schlein~\cite{BBCOS}, together with intermediate operator estimates
of Fournais, Junge, Girardot, Morin, Olivieri, and Triay~\cite{FJGMOT}.
Their positive excitation terms give quantitative closeness to the
local quadratic vacuum. We extract from them a uniform estimate for
one-body operators supported on the local excitations, including operators
whose matrices are complex and not diagonal in momentum. This stronger
observable estimate will allow us to test translations after the
long-wavelength contribution has been controlled.

The second term concerns wavelengths larger than $\ell$. To control it,
we first take a normalized nonnegative ground vector $\Phi$. For a cube $v$ in a
torus of volume $V$, define a measure on the remaining $N-1$ coordinates
$Y'$ by the density
\begin{equation}\label{intro:palm}
 s_v(Y')^2=\frac{V}{|v|}\int_v\Phi(y,Y')^2\,\dd y.
\end{equation}
This is a finite smoothed particle-deletion measure: the deleted
particle is selected with weight proportional to its presence in $v$,
and its position is integrated out. Its mass is the expected particle
count in $v$ divided by $(N/V)|v|$; this mass averages to one over
equal-volume cells.
Comparing the square roots $s_v$ for distant
cubes gives control of spatial variation that the energy cannot supply.

The comparison transports the deletion along a chain of particles.
Two copies of the configuration are coupled so that deleting the first
particle in one copy and the last in the other leaves identical remaining
configurations.
Short Brownian bridges implement the rearrangement while preserving
the exact hard-sphere constraints. The normalization of the conditional
bridge law cancels in the likelihood ratio. Near either end of a long
chain, the paths in one copy are kept unchanged. These unchanged segments
remove endpoint interactions from the errors measured by a second moment.
Random routes with sparse intersections then make that error smaller
than the depletion scale.

Three estimates connect these deletion measures to the projection.
A fourth-moment comparison between neighboring small-cell root mean
squares shows that their arithmetic mean is close to their root mean
square over a cell of side $\ell$.
The occupation estimate on those same small cells compares their root
mean squares with arithmetic averages of $\Phi$. A second-moment
comparison between distant cells of side $\ell$ makes their root mean
squares nearly constant. Together these estimates bound the last term
of \eqref{intro:split} by
$o(\sqrt{\rho a^3})$. A decomposition into real and imaginary positive
parts extends that bound to every ground vector.

It remains to identify the momentum distribution of these local
excitations. Translation by a vector $h$ multiplies a plane wave of
momentum $k$ by $e^{ik\cdot h}$, so expectations of translations are
Fourier transforms of the occupation measure. We restrict a translation
to large boxes and average over shifts of the box partition. Only a small
fraction of the shifted partitions separate a point from its translate.
Sandwiching the translation between projections onto local excitations
makes this error proportional to the depleted occupation. The local
one-body estimate then gives the Bogoliubov Fourier transform and
convergence on compactly supported tests. Together with total depletion,
this supplies tightness and convergence against every bounded continuous
function. All estimates are
uniform in pure ground vectors, and convexity gives the mixed-state
assertion.

\paragraph{Organization.}
Section~\ref{sec:setup} sets the scaling and states the probabilistic
inputs. Section~\ref{sec:energy} retains the positive excitation operators
in the local energy bound, and Section~\ref{en:consequences} extracts
the count and one-body estimates.
Sections~\ref{sec:local-paths} and~\ref{sec:local-palm}
construct the short-path comparison and prove its fourth-moment bound.
Sections~\ref{sec:routes} and~\ref{sec:long-palm} construct the long routes
and compare distant deletion measures. Section~\ref{sec:assembly}
combines the estimates to control the long wavelengths.
Section~\ref{sec:momentum} uses averaged translations to prove
Theorem~\ref{thm:main}.
